\subsection{What purity gives in a differential module}\label{sec:seedbridge}

The following lemma is the precise passage from the formal Gaussian
functional to an analytic period. Its hypotheses must be verified
after each change of variables or localization.

\begin{lemma}[The seed functional]\label{lem:seedfunctional}
Let $M$ be a finite-dimensional differential module over $K=\C(t)$,
let $s\in M$, and suppose a $K$-linear horizontal map
$J:M\longrightarrow\C((t))$ satisfies $J(s)=1$. Then the differential
submodule $U=K\langle\partial_t\rangle s$ has a horizontal quotient
$L:U\longrightarrow K$ with $L(s)=1$.

At an ordinary base point, $L$ extends locally to an analytic horizontal
functional on $M$. For the period conclusion, assume explicitly that
$M$ is the generic exponential de Rham system of regular functions
$f,g$ on a smooth affine complex variety, where $f$ is the base
projection and $f,g$ are algebraically independent. Remove the finite
exceptional base set in the Hien--Roucairol comparison theorem.
Then the extension is represented by a flat rapid-decay cycle, and its
period evaluated on $s$ is identically one.
\end{lemma}
\begin{proof}
On the generators $s,\nabla s,\nabla^2s,\ldots$, the values of $J$
are $1,0,0,\ldots$. Therefore its restriction takes values in $K$.
Every rational relation between these generators is respected because
$J$ was defined on $M$, so this restriction is well defined and
horizontal. Its value at $s$ makes it surjective.

Choose a point where $M$, $U$ and $L$ are ordinary, extend the fiber
functional from $U$ to $M$, and solve the dual analytic ODE.
Uniqueness preserves its restriction to $U$. Perfectness of the period pairing is
Hien--Roucairol~\cite[Theorem~2.7]{HienRoucairol}; their
Theorem~3.5 represents holomorphic solutions by rapid-decay periods.
Together these statements give the asserted cycle. A scalar rational normalization is an invertible gauge
away from its zeros and poles and must be retained in the period.
\end{proof}

There are three distinct verifications here. First, integration by
parts must make $J$ descend to the exact displayed de Rham quotient.
Second, every inverted spatial coordinate must be a valid formal
Laurent unit on the selected branch. Third, the comparison theorem
requires a smooth affine total space and regular functions after
removing the finite set of poles in the parameter. The theorem gives
a representation, not a bound at infinity. In particular, it gives
no shrinking estimate on its own.

There is also a useful purely formal way to obtain a contradiction.

\begin{lemma}[A full formal frame and a positive seed]\label{lem:positiveformal}
Assume the hypotheses of the first paragraph of
Lemma~\ref{lem:seedfunctional}. Suppose that, after a finite ramification
$r=t^{-1/N}$ at infinity, $M$ has a full formal frame of rank-one
exponential factors. For every factor with zero exponential part,
assume its horizontal functional takes $s$ into $r\C[[r]]$.
Then $J(s)=1$ is impossible.
\end{lemma}
\begin{proof}
Enlarge the ramification to include the rational powers in the frame.
The module over $\C((r))$ is then a direct sum of rank-one modules,
so the quotient on $U$ extends to a horizontal map from the whole
module to $\C((r))$. A nonzero exponential factor has no such map:
the logarithmic derivative of a Laurent series has at most a simple
logarithmic pole and cannot equal the derivative of a nonzero
negative-power exponential part. On the zero part, all horizontal
functionals are constant linear combinations of the supplied frame.
Their values on $s$ have strictly positive $r$-order and cannot sum
to one.
\end{proof}

The word \emph{full} is essential. Counting some critical points or
writing a few scalar asymptotic expansions does not supply this lemma's
hypothesis. The finite and escaping pieces must together give an
invertible period matrix of the full de Rham rank. Also, a fractional
power becomes an integer after ramification; the proof uses its
strictly positive order, not an unsupported monodromy inference.
